\documentclass[10pt,a4paper]{amsart}
\usepackage[margin=19mm]{geometry}
\usepackage{amsmath,amssymb,mathtools}
\usepackage[scr=rsfs,cal=euler]{mathalfa}
\usepackage{xfrac}
\usepackage[unicode,hidelinks,bookmarks=false]{hyperref}
\newcommand{\ie}{i.e.\ }
\newcommand{\cf}{cf.\ }
\newcommand{\resp}{respectively\ }
\newcommand{\Subsection}{\subsection}
\providecommand{\nobreakdash}{\nobreak\hbox{-}}
\setlength{\emergencystretch}{2em}
\multlinegap0pt
\usepackage{bm}
\usepackage{bbm}
\usepackage{dsfont}
\usepackage{xspace}
\usepackage{cancel}
\usepackage{mathtools}
\usepackage{amsthm}
\makeatletter
\@ifundefined{lemma}{\newtheorem{lemma}{Lemma}}{}
\@ifundefined{proposition}{\newtheorem{proposition}{Proposition}}{}
\@ifundefined{theorem}{\newtheorem{theorem}{Theorem}}{}
\makeatother
\title[The Antisymmetric Line Graph]{The Antisymmetric Line Graph}
\author{Hartosh Singh Bal}
\date{Source: arXiv:2603.03087v2 (CC BY 4.0)}
\begin{document}
\maketitle

\noindent\emph{Source and licence.} The Antisymmetric Line Graph. Version arXiv:2603.03087v2 (CC BY 4.0), distributed under the Creative Commons Attribution 4.0 International licence, as recorded in the arXiv licence metadata for this exact version and shown on the abstract page https://arxiv.org/abs/2603.03087v2. Full rights notice: licenses/arxiv-2603.03087-CC-BY-4.0.txt.\par\smallskip

\noindent\textit{Editorial scope.} Counted unit: the Proposition whose source label is \texttt{prop:complete-multipartite} in the article (source0.tex lines 1852--1874, source label \texttt{prop:complete-multipartite}), delivered together with the article's own complete original proof of it (source0.tex lines 1876--1955, one \texttt{proof} environment of 80 lines). No mathematics is written, repaired or re-derived: statement and proof are the article's own text line for line. Required context retained uncounted: lem:orientation-frustration (lines 1166--1172); thm:spectral-lower (lines 1655--1668). Every definition, display and label of those lines is retained unchanged. Omitted from this package and not redistributed: 1--1165, 1173--1654, 1669--1851, 1875--1875, 1956--2166. Nothing inside a retained excerpt is omitted or altered: every retained body is delivered exactly as the article's lines are written, so no omission receipt of any kind accompanies this package. Citations: the retained excerpts carry no citation command at all, so no bibliography entry of the article is transcribed and no reference is redistributed. Reference adaptation: none was needed and none was made; every reference inside the delivered unit resolves inside it, so no unresolved reference or citation is produced. Printed slips kept literal: no printed word, symbol, hypothesis or number of the counted statement or of its proof was altered. Layout and class: the article's own class is \texttt{article} with packages including amsmath, amsthm, amssymb, graphicx, geometry, amsfonts, hyperref, tikz, booktabs, xy, braket, algorithm, algorithmic, subcaption; the wrapper is the lane's pinned \texttt{amsart} editorial wrapper at 10pt on A4, which restates the theorem-like environments the retained text needs and adds the article's own macro definitions that the retained text needs (none), with extra wrapper packages (bm, bbm, dsfont, xspace, cancel, mathtools). The delivered numbering is the unit's own. Assets: no retained span contains a figure environment, a graphics-inclusion command, a PDF-page inclusion, a table or a TikZ picture; no image file is copied into this package, the package declares no asset dependency and no image file is redistributed with it. Licence: version arXiv:2603.03087v2 of this article is distributed under the Creative Commons Attribution 4.0 International licence, as recorded in the arXiv licence metadata for this exact version and shown on the abstract page https://arxiv.org/abs/2603.03087v2. A full rights notice is delivered at licenses/arxiv-2603.03087-CC-BY-4.0.txt. Characters: every retained excerpt is pure ASCII, so the delivered engine, XeLaTeX with its default Latin Modern text font, reports no missing character.
\begin{lemma}[Orientation characterization]\label{lem:orientation-frustration}
For any finite simple graph $G$, one has
\[
\ell(\mathcal A(G)) \;=\; \min_{\vec{G}} \;\sum_{v\in V(G)} d^+_{\vec G}(v)\,d^-_{\vec G}(v),
\]
where the minimum is taken over all reference orientations $\vec G$ of $G$.
\end{lemma}

\begin{theorem}[Spectral lower bound]
\label{thm:spectral-lower}
Let
\[
0=\lambda_1\le \lambda_2\le \cdots \le \lambda_n
\]
be the Laplacian eigenvalues of $G$, so that $\lambda_n=\lambda_{\max}(L)$.
Then
\[
\ell(\mathcal A(G))
\ge
\frac{\sum_{v\in V(G)} d(v)^2-|E(G)|\,\lambda_n}{4}.
\]
\end{theorem}

\begin{proposition}[Complete multipartite graphs: exact evaluation and $3/4$ sharpness]
\label{prop:complete-multipartite}
Let
\[
G=K_{n_1,\dots,n_q}
\]
be the complete multipartite graph with part sizes $n_1,\dots,n_q$, and write
\[
n=\sum_{i=1}^q n_i.
\]
Then
\[
\ell(\mathcal A(G))
=
\sum_{1\le i<j<k\le q} n_i n_j n_k.
\]
Moreover, the spectral lower bound from Theorem~\ref{thm:spectral-lower} is
\[
\ell(\mathcal A(G))
\ge
\frac34\sum_{1\le i<j<k\le q} n_i n_j n_k.
\]
\end{proposition}

\begin{proof}
By Lemma~\ref{lem:orientation-frustration},
\[
\ell(\mathcal A(G))
=
\min_{\vec G}\sum_{v\in V(G)} d^+_{\vec G}(v)d^-_{\vec G}(v),
\]
and the right-hand side counts directed $2$--paths.

Every triangle of $G$ is obtained by choosing one vertex from each of three
distinct parts. Hence
\[
\#\{\text{triangles of }G\}
=
\sum_{1\le i<j<k\le q} n_i n_j n_k.
\]
In any orientation of a triangle, there is at least one directed $2$--path.
Therefore every orientation of $G$ has at least one directed $2$--path per
triangle, so
\[
\ell(\mathcal A(G))
\ge
\sum_{1\le i<j<k\le q} n_i n_j n_k.
\]

For the reverse inequality, order the parts
\[
V_1,\dots,V_q
\]
and orient every edge from the earlier part to the later part.
Then every triangle is transitive and contributes exactly one directed
$2$--path.
At a vertex in part $V_j$, all incoming edges come from parts $V_i$ with $i<j$
and all outgoing edges go to parts $V_k$ with $k>j$, so every directed
$2$--path through that vertex has endpoints in two distinct parts and hence lies
in a unique triangle.
Thus the total number of directed $2$--paths is exactly the number of triangles,
and therefore
\[
\ell(\mathcal A(G))
=
\sum_{1\le i<j<k\le q} n_i n_j n_k.
\]

It remains to evaluate the spectral lower bound.
For a complete multipartite graph one has $\lambda_n(G)=n$.
Also, every vertex in part $i$ has degree $n-n_i$, so
\[
\sum_{v\in V(G)} d(v)^2=\sum_{i=1}^q n_i(n-n_i)^2,
\]
and
\[
|E(G)|=\frac{n^2-\sum_{i=1}^q n_i^2}{2}.
\]
Hence Theorem~\ref{thm:spectral-lower} gives
\[
\ell(\mathcal A(G))
\ge
\frac{1}{4}\left(\sum_{i=1}^q n_i(n-n_i)^2
-\frac{n}{2}\left(n^2-\sum_{i=1}^q n_i^2\right)\right).
\]
Expanding and simplifying yields
\[
\ell(\mathcal A(G))
\ge
\frac18\left(n^3-3n\sum_{i=1}^q n_i^2+2\sum_{i=1}^q n_i^3\right).
\]
On the other hand, the elementary symmetric identity
\[
\sum_{1\le i<j<k\le q} n_i n_j n_k
=
\frac16\left(n^3-3n\sum_{i=1}^q n_i^2+2\sum_{i=1}^q n_i^3\right)
\]
shows that the spectral lower bound equals
\[
\frac34\sum_{1\le i<j<k\le q} n_i n_j n_k
=
\frac34\,\ell(\mathcal A(G)).
\]
\end{proof}

\end{document}
